%% ==========================================================================
%%  FINAL REPORT --- SECTION 7, LEAKAGE OF THE COMPOSITION
%%  Sources: tech_leakage.tex (mid-term, rewritten), docs/threat-model.md
%%  sections 4 and 9, model/attack.py, model/dp_study.py,
%%  model/poison_study.py, bench/bench_dos.cpp, tests/test_harvest.cpp.
%%  Figures F1 F6 F10 F11 live here, each exactly once.
%% ==========================================================================
\section{Leakage of the composition}
\label{sec:leak}

Section~\ref{sec:eval} priced the private read. This section asks what it buys:
what an observer learns when the fetch is \emph{not} private, whether anything
cheaper than a private read closes the gap, and what a malicious client can do
to the composed system. Neither base paper can answer these, because neither
implements both halves.

\subsection{Why a fetch is a projection}
\label{sec:leak:why}

$\itemembed$ is published, so every observer holds a complete latent-factor model
of the catalogue. The items recommended to user $i$ are the ones whose columns of
$\itemembed$ score highest against the private embedding $\avec{i}$. An observer
who sees \emph{which} items the user fetched therefore learns $\topk$ vectors
that all point roughly where $\avec{i}$ points. Each fetch is a constraint on the
private embedding, expressed in a basis the observer already has. The question is
how quickly those constraints give it away.

\subsection{The attack, and the control that makes it a result}
\label{sec:leak:attack}

The adversary is given $\itemembed$ and the indices of the first $j$ records a
user fetched, and nothing else: no scores, no ratings, no part of $\userembed$.
It estimates $\hat{\mathbf{a}}$ as the mean of the fetched items' unit-normalised
embeddings. We report the cosine between $\hat{\mathbf{a}}$ and the true
$\avec{i}$, and the share of the user's \emph{next} twenty recommendations that
the ranking induced by $\hat{\mathbf{a}}$ predicts, excluding the $j$ items
already seen. All 943 users are attacked, against the same model the demo serves.

\paragraph{The control.} Popular items sit near the top of almost everyone's
ranking, so an estimator can look strong while having learned only what is
public. The control therefore runs the identical estimator on the $j$ globally
most popular items, ignoring the user entirely, and the attack may claim only what
it achieves above that line. A random direction is also reported, to show where
the floor of pure guessing lies.

\begin{table}[t]
\centering
\small
\caption{Recovering a user's private embedding from observed fetches alone,
averaged over all 943 users. Data: \texttt{bench/results/leakage\_attack.jsonl}.}
\label{tab:attack}
\begin{tabular}{@{}rcccc@{}}
\toprule
Fetches $j$ & $\cos(\hat{\mathbf{a}}, \mathbf{a})$ & popularity control &
  margin & next-20 overlap \\
\midrule
1  & 0.55 & 0.26 & 0.28 & 39\% \\
2  & 0.66 & 0.26 & 0.40 & 47\% \\
5  & 0.77 & 0.39 & 0.38 & 53\% \\
10 & \textbf{0.84} & 0.44 & 0.40 & \textbf{56\%} \\
20 & 0.89 & 0.47 & 0.42 & 55\% \\
50 & 0.91 & 0.53 & 0.39 & 46\% \\
\bottomrule
\end{tabular}
\end{table}

\begin{figure}[t]
\centering
\includegraphics[width=\linewidth]{fig_attack}
\caption{Left: cosine between the estimate and the true embedding, with the
interquartile band of the attack. Right: the share of the user's next twenty
recommendations the adversary names. The gap between the attack and the
popularity control is the part that is about the individual. The second attack
curve is a stronger-looking estimator that uses the ranking constraints, and it
does worse. Data: \texttt{bench/results/leakage\_attack.jsonl}.}
\label{fig:attack}
\end{figure}

\textbf{One observed fetch gives a cosine of 0.55, twice the popularity control.
Ten give 0.84, enough to name 56\% of the user's next twenty recommendations}
(Table~\ref{tab:attack}, Figure~\ref{fig:attack}). The margin over the control is
0.28 at one fetch and between 0.38 and 0.42 from two fetches on. The random
direction averages 0.003.

We raise three caveats ourselves. First, one user's cosine is not evidence. The
random control's interquartile range runs from $-0.17$ to $0.18$, which is what a
random direction in sixteen dimensions gives, so a single user at 0.2 is inside
the noise of guessing. Only population means and the margin over the control
carry a claim. Second, the overlap column rises and then falls, peaking near ten
fetches. That is the metric, not the attack: the observed items are removed from
the candidates, so the easy head of the ranking is progressively taken away,
while the cosine keeps rising. Third, a second estimator that uses the fact that
every fetched item outranks every unfetched one, strictly more information than
the mean uses, does \emph{worse} from two fetches on, 0.81 against 0.91 at fifty.
That sharpens the conclusion. The cheap, obvious attack is already enough, and an
adversary needs no sophistication.

\subsection{Defence one: decoys}
\label{sec:leak:decoys}

The obvious cheap alternative to a private read is to fetch the $\topk$ real
recommendations together with $r$ decoys. The decoys must be drawn in proportion
to popularity, because real recommendations are popular and uniformly drawn
decoys would come from the long tail, where an observer could discard them by
inspection and attack what remained.

\begin{figure}[t]
\centering
\includegraphics[width=\linewidth]{fig_decoy}
\caption{The attack at ten real fetches plus $r$ popularity-weighted decoys,
against decoys per query (left) and against traffic per query (right). The
dashed line is the popularity control. Data:
\texttt{bench/results/leakage\_attack.jsonl}.}
\label{fig:decoy}
\end{figure}

It does not work (Figure~\ref{fig:decoy}). Forty decoys, five times the real
traffic, lower the attack's cosine only from 0.843 to 0.705, against a
popularity control at 0.442. The reason is the same fact that made the control
necessary. Popularity-weighted decoys point along the popularity direction, which
is exactly what the control already recovers, so decoys drag the estimate
\emph{towards} the control and no further. The defence has a floor built into
it, and the floor is where an adversary who watched nobody already stands.

\subsection{Defence two: \nudgeplain{}'s differential privacy}
\label{sec:leak:dp}

\nudgeplain{}'s Section~9 extends the system with differential privacy. The
servers normalise each user's row of $\ratings$, share a symmetric Gaussian
matrix $E$ once through their PRF keys, and add $E\mathbf{v}$ at every step of
power iteration, so the iteration runs on $\ratings^{\top}\ratings + E$ at no extra
communication. The paper does not state the noise scale. We take it from the
Gaussian mechanism of Dwork, Talwar, Thakurta and Zhang that it
cites~\cite{analyzegauss2014}: with every row of unit norm, one user changes
$\ratings^{\top}\ratings$ by at most one in Frobenius norm, and $\sigma =
\sqrt{2 \ln (1.25/\delta)}/\varepsilon$ gives $(\varepsilon, \delta)$-differential
privacy~\cite{dwork2014}. We use $\delta = 2^{-40}$, as \nudgeplain{} does.

We deviate from the paper in one place, deliberately. \nudgeplain{} normalises
row $i$ by dividing by the user's number of ratings $c_i$. With ratings up to
five, that bounds the row's norm by $5/\sqrt{c_i}$, which exceeds one whenever
$c_i < 25$, and in the training split we use, 200 of the 943 users have fewer than
25 ratings. For them the privacy bound would not hold, so we normalise each row to
unit $\ell_2$ norm instead. The count-normalised arm was run too and is in the data. It behaves the
same way.

\begin{table}[t]
\centering
\small
\caption{\nudgeplain{}'s mechanism on MovieLens-100K at $\delta = 2^{-40}$,
with the attack at ten fetches. Data: \texttt{bench/results/dp\_study.jsonl}.}
\label{tab:dp}
\begin{tabular}{@{}lccc@{}}
\toprule
 & nDCG@20 & attack cosine & margin over control \\
\midrule
no noise (rows normalised) & 0.4255 & 0.8231 & 0.4324 \\
$\varepsilon = 8$ & 0.1361 & 0.9167 & 0.5044 \\
$\varepsilon = 1$ & \textbf{0.1239} & \textbf{0.9212} & \textbf{0.4898} \\
\bottomrule
\end{tabular}
\end{table}

\paragraph{Both axes move the wrong way.} At $\varepsilon = 1$, quality falls by
71\% and the attack gets \emph{stronger} (Table~\ref{tab:dp}). Even at
$\varepsilon = 8$ quality falls to 0.136. The second result is not a paradox.
Differential privacy here protects the \emph{training data}: it bounds what
$\itemembed$ reveals about any user's ratings. The attack recovers \emph{this}
user's embedding using $\itemembed$ as a basis, and the user and the adversary
work against the same published, noisy matrix, so the noise does not disturb the
attack's geometry. What the noise does destroy is the shared popularity structure
the control exploits, which widens the attack's margin.

\paragraph{Why we lose 71\% where \nudgeplain{} loses 17\%.} \nudgeplain{}
reports nDCG@20 falling from 0.29 to 0.24 on Netflix~\cite{nudge2026}. The
mechanism is the same, and the dataset is not. The noise matrix's spectral norm
grows as $\sigma\sqrt{\items}$ while the signal, the spectral norm of
$\ratings^{\top}\ratings$, grows with the number of users $\users$. The usable
regime is set by $\users/\sqrt{\items}$, which is 157 times larger for Netflix,
with 480{,}189 users and 17{,}770 items, than for MovieLens-100K. Measured
directly, $\norm{\ratings^{\top}\ratings}_2 = 144.6$ against $\norm{E}_2 \approx
613$ at $\varepsilon = 1$, so the noise dominates the signal by a factor of 4.2
(Figure~\ref{fig:dp}). This is a statement about our dataset's scale, not a
contradiction of the paper.

\begin{figure}[t]
\centering
\includegraphics[width=\linewidth]{fig_dp}
\caption{Left: both cheaper defences in one plane, attack strength against
recommendation quality. A defence that worked would reach the popularity
control. Right: the spectral signal-to-noise ratio against $\varepsilon$ for our
dataset and for Netflix's dimensions. Data:
\texttt{bench/results/dp\_study.jsonl} and
\texttt{bench/results/leakage\_attack.jsonl}.}
\label{fig:dp}
\end{figure}

Neither cheaper defence reaches the popularity control. Decoys cost bandwidth and
stay far above it, and differential privacy costs most of the model's quality and
moves the attack further away from it. That is the argument for making the read
itself private, and it is the gap \nudgeplain{} leaves to ``other means''.

\subsection{The malicious client}
\label{sec:leak:malicious}

The last adversary class the composition creates is a client who does not follow
the protocol. For retrieval alone, a malformed query only corrupts that client's
own answer. The harvest changes that, because a query now also writes into the
next model.

\paragraph{The hole.} An honest client generates its keys with $\beta = 1$, so
the two expansions differ by exactly the indicator of the fetched item. Nothing
stopped a client from generating keys with any other $\beta$. Retrieval still
works, since the record comes back scaled by $\beta$ and the client divides it
out, and the servers add $\beta$, not one, into the consumption counts. Every
existing check passes: the keys parse, the domain matches, the expansion has the
right length. Our test shows one query at $\beta = 10^6$ placing 1{,}000{,}000
into the next round's training input.

\paragraph{What it buys.} Less than we first assumed, and for an instructive
reason (Figure~\ref{fig:malicious}, left). A single attacker at weight $10^6$
lowers honest users' nDCG@20 by between 0.7 and 1.0\%, depending on its targets,
and the items it promotes do not rise. Their mean rank worsens, from 901 to 973
of 1{,}682. Power iteration normalises at every step,
so beyond some point extra weight only fixes a direction, and one attacker
captures one of the $\dime = 16$ components while the other fifteen still carry
the honest signal. Colluders with distinct targets are another matter. Eight cost
4.1\%, and sixteen take nDCG@20 from 0.4536 to 0.0021. We ran only powers of two,
so the data place the collapse somewhere between eight and sixteen colluders,
and it is consistent with the attackers capturing all $\dime$ components. This is
destruction of the model, not promotion of items.

\paragraph{The check.} Summation is linear, so each PIR server can sum its own
expansion locally, and the two sums differ by exactly $\beta$, whatever the index.
The servers open that one scalar and require it to equal one: one round and two
ring elements per query. Opening it reveals nothing about the index, because
$\beta$ is meant to be the public constant one and each server's own sum is
pseudorandom. Our test confirms the check accepts $\beta = 1$ and rejects
$\beta = 10^6$ at that cost.

\paragraph{What the check does not close.} It bounds the total weight of a query,
not how the weight is distributed. A client that forges the expansion directly,
rather than calling the key generator, can place $+2$ at one item and $-1$ at
another, which sums to one and passes. Our test constructs exactly that forgery
and shows it passes. A forgery of the same shape with large opposite weights would
give colluders the same power as before, and we have not measured it. Closing it
requires proving that a key encodes a genuine point function, which is what the
audit of \textsc{Sabre}~\cite{sabre2022} does, and we did not build it.

\paragraph{Denial of service.} Answering any query means expanding the key over
the whole domain, while the key is only logarithmic in the domain. So a query
buys server work out of proportion to its size (Figure~\ref{fig:malicious},
right). At our 11-bit domain, a 227-byte key buys 142 times its own cost in
server time in the run the figure shows, and 121 times in an earlier committed
run. At 16 bits the same run measures 4{,}323 times. The amplification grows with
the catalogue, the wrong direction for a defender. Malformed keys are cheap to reject,
since parsing validates them before any expansion, so the expensive attack is the
one that looks legitimate. The linear work is not an inefficiency. It is the
privacy property, since any shortcut that depended on the index would reveal it.
This is the clearest example in the project of privacy paid for in availability,
and the same unbuilt audit is what would address it.

\begin{figure}[t]
\centering
\includegraphics[width=\linewidth]{fig_malicious}
\caption{Left: honest users' nDCG@20 against the number of colluding clients,
each voting at weight $10^6$ for distinct targets, before the weight check. Only
powers of two were run. Right: server work bought by one key, and the key's size,
against the domain. The star is our operating point. Data:
\texttt{bench/results/poison\_study.jsonl} and
\texttt{bench/results/bench\_dos.jsonl}.}
\label{fig:malicious}
\end{figure}
